\documentclass{beamer}

\usepackage{kantlipsum}

\begin{document}

	\begin{frame}[allowframebreaks, t]
		\frametitle{Title of the Frame}
		\framesubtitle{Subtitle for the Frame}
		\kant[1]
		The use of this option is evil. In a (good) presentation you prepare each slide carefully and think
twice before putting something on a certain slide rather than on some different slide. Using the
allowframebreaks option invites the creation of horrible, endless presentations that resemble more
a “paper projected on the wall” than a presentation. Nevertheless, the option does have its uses.
Most noticeably, it can be convenient for automatically splitting bibliographies or long equations.
	\end{frame}
	
	\begin{frame}[fragile, t]{Frame with Verbatim Text}
		\begin{verbatim}
			\begin{LaTeX}
				% Some \LaTeX Code Here !!
			\end{LaTeX}
		\end{verbatim}
	\end{frame}
	
	\begin{frame}[allowframebreaks, allowdisplaybreaks]
		\begin{align*}
			x&=y \\ x&=y \\ x&=y \\ x&=y \\ x&=y \\ 
			x&=y \\ x&=y \\ x&=y \\ x&=y \\ x&=y \\ 
			x&=y \\ x&=y \\ x&=y \\ x&=y \\ 
			x&=2*y \\ 
			x&=y \\ x&=y \\ x&=y \\ x&=y \\ x&=y 
		\end{align*}
	\end{frame}
	
\end{document}